\section{Expanded Limitations}
\label{app:limitations}

\noindent\textbf{The working posterior is not identified by a probability map.}
The shared-threshold construction selects one coupling from many with the same pixel
marginals. Marginal calibration, even if exact, does not identify expected nonlinear
geometric loss. Repeated annotations or another source of label variability would be
needed to validate this coupling directly; one reference mask per image cannot
separate marginal calibration error from dependence-model error. Consequently,
\(W_{1,d}(P_x,Q_{p(x)})\) is generally unobservable in the present datasets:
the transport bounds structure model error but are not computable deployment
certificates. They can be evaluated directly only with a known synthetic
posterior or sufficiently rich repeated annotations.

\noindent\textbf{Geometry compatibility does not validate the posterior.}
The Dice--Jaccard and HD--surface results show that their respective costs are
Lipschitz in the stated ground geometries. They do not show that
\(Q_{p(x)}\) is close to \(P_x\) in either geometry. Likewise, Exact Dice
removes quadrature error only after \(Q_p\) has been chosen. Mask-space total
variation is a formally valid fallback but can be nearly one when two discrete
laws have different support even if their masks are geometrically close.

\noindent\textbf{The loss convention is substantive.}
Dice is naturally bounded, but surface distances require an explicit empty-mask
rule. Assigning a one-sided empty comparison the image diagonal, a fixed cap, or an
application-specific cost produces different losses and can change the ranking. Our
penalized HD and HD95 in unit-diagonal coordinates are two fixed choices. The published SDC baseline's
empty--empty Dice value also differs from ours, so its theoretical approximation
guarantee does not automatically cover empty deployed masks under our loss.
Dropping undefined cases would change the target population and is not equivalent.
The SDC count-dispersion bound also scales as \(1/|\Yhat|\), so it can be loose
for small predicted objects or diffuse foreground probability outside the
action; it explains when SDC is close to expected Dice rather than guaranteeing
uniform closeness.

\noindent\textbf{Quadrature control is not posterior control.}
\Cref{eq:decomp} separates numerical and posterior terms but does not make
the latter small by assertion, and the variation constant can be loose. The
population \(L^1\) bound yields the AURC-regret envelope in
\Cref{app:cor:aurc-regret}, but its posterior term remains even as
\(M\to\infty\), and the resulting control is weakest at very low coverage. Thus a
larger \(M\) need not improve empirical AURC monotonically. A pairwise margin is
still needed to guarantee exact ordering; aggregate AURC control does not provide
rank recovery for every close pair.

\noindent\textbf{Scope and information limits.}
We focus on native binary masks and single-map, post-hoc confidence. If two cases
produce the same \((p,\Yhat)\), every deterministic score studied here gives them
the same confidence even when external context makes one more likely to fail.
Features, ensembles, auxiliary failure predictors, distribution-shift detectors,
multiclass actions, and multi-rater posterior validation are outside the present
scope. Finally, AURC measures ranking; it does not imply that a confidence value is
calibrated as a probability or an absolute forecast of loss.

\noindent\textbf{Empirical generality and training variability.}
Even the fixed five-dataset primary design and the separately locked
SegFormer/DUTS extension
cannot establish that risk-aligned
confidence will dominate across architectures, shifts, or annotation protocols.
Our paired comparisons deliberately condition on each fixed probability map and
therefore isolate confidence ranking, not segmentation-model accuracy. The paired
image bootstrap quantifies held-out image variation for a fixed fitted
checkpoint; its percentile intervals are pointwise and do not carry a
simultaneous family-wise-error guarantee. The primary target-model analysis and
its bootstrap intervals condition on seed 0. A separate extension repeats the
same protocol at seeds 1 and 2 only for the ten target conditions and summarizes
checkpoint-level variation descriptively. The prespecified training-sensitivity
rule is met specifically for HD versus HD95 under HD95 risk: five of ten target conditions reverse
direction across the three checkpoints, and seed 0 is outside the majority direction
in three. We therefore call only this comparison training-sensitive. These three
checkpoints neither estimate the full distribution of training randomness nor
support inference over optimization or training-sample variation, and they do not
enter the primary intervals. Repeated training runs at larger scale and external-site
validation would be needed for that broader claim. Moreover, model
conditions reuse the same held-out images within each dataset;
condition-specific rows are paired evaluations, not independent images.
Release-level patient, video, or eye identifiers are not available uniformly
across the five cohorts, so the intervals should not be read as patient-clustered
evidence.
